\documentclass[11pt]{article}
\usepackage[T1]{fontenc}
\usepackage[utf8]{inputenc}
\usepackage{lmodern}
\usepackage[letterpaper,margin=0.78in]{geometry}
\usepackage{microtype}
\usepackage{booktabs,longtable,tabularx,array,multirow}
\usepackage{amsmath,amssymb}
\usepackage{xcolor}
\usepackage{enumitem}
\usepackage{titlesec}
\usepackage{fancyhdr}
\usepackage{lastpage}
\usepackage{hyperref}
\usepackage[most]{tcolorbox}
\usepackage{parskip}
\usepackage{needspace}
\usepackage{ragged2e}
\usepackage{setspace}

\definecolor{Navy}{HTML}{17365D}
\definecolor{Blue}{HTML}{2F75B5}
\definecolor{Teal}{HTML}{1F6D6A}
\definecolor{Gold}{HTML}{B38B2D}
\definecolor{Gray}{HTML}{5E6B75}
\definecolor{LightBlue}{HTML}{EEF4FA}
\definecolor{LightGold}{HTML}{FBF6E8}
\definecolor{LightGray}{HTML}{F4F5F7}
\definecolor{LightTeal}{HTML}{EDF7F6}

\hypersetup{colorlinks=true,linkcolor=Navy,urlcolor=Blue,citecolor=Teal,pdfauthor={OpenAI},pdfsubject={Value investing study and implementation materials}}
\setlist[itemize]{leftmargin=1.5em,itemsep=2pt,topsep=3pt}
\setlist[enumerate]{leftmargin=1.7em,itemsep=2pt,topsep=3pt}
\renewcommand{\arraystretch}{1.16}
\setlength{\tabcolsep}{5pt}
\titleformat{\section}{\Large\bfseries\color{Navy}}{\thesection}{0.6em}{}[\vspace{-0.25em}\color{Blue}\rule{\linewidth}{0.5pt}]
\titleformat{\subsection}{\large\bfseries\color{Teal}}{\thesubsection}{0.5em}{}
\titleformat{\subsubsection}{\normalsize\bfseries\color{Gray}}{\thesubsubsection}{0.5em}{}
\pagestyle{fancy}
\fancyhf{}
\fancyhead[L]{\small\color{Gray}\DocShortTitle}
\fancyhead[R]{\small\color{Gray}Value Investing Reference Set}
\fancyfoot[C]{\small\color{Gray}Page \thepage\ of \pageref*{LastPage}}
\renewcommand{\headrulewidth}{0.4pt}
\renewcommand{\footrulewidth}{0pt}

\newtcolorbox{principlebox}[1][]{enhanced,breakable,colback=LightBlue,colframe=Blue,boxrule=0.6pt,arc=2mm,left=2mm,right=2mm,top=1.5mm,bottom=1.5mm,title={#1},fonttitle=\bfseries}
\newtcolorbox{cautionbox}[1][]{enhanced,breakable,colback=LightGold,colframe=Gold,boxrule=0.6pt,arc=2mm,left=2mm,right=2mm,top=1.5mm,bottom=1.5mm,title={#1},fonttitle=\bfseries}
\newtcolorbox{modernbox}[1][]{enhanced,breakable,colback=LightTeal,colframe=Teal,boxrule=0.6pt,arc=2mm,left=2mm,right=2mm,top=1.5mm,bottom=1.5mm,title={#1},fonttitle=\bfseries}
\newtcolorbox{worksheetbox}[1][]{enhanced,breakable,colback=white,colframe=Gray,boxrule=0.45pt,arc=1mm,left=2mm,right=2mm,top=1.5mm,bottom=1.5mm,title={#1},fonttitle=\bfseries}
\newcommand{\checkbox}{\(\square\)}
\newcommand{\blankline}{\rule{0.96\linewidth}{0.35pt}}
\newcommand{\sourceRef}[1]{\textit{Source anchor: #1}}
\newcolumntype{Y}{>{\RaggedRight\arraybackslash}X}
\newcolumntype{C}[1]{>{\centering\arraybackslash}p{#1}}

\newcommand{\makedoccover}[4]{%
\begin{titlepage}
\thispagestyle{empty}
\vspace*{1.0in}
{\color{Navy}\rule{\linewidth}{2.2pt}}\\[0.55in]
{\Huge\bfseries\color{Navy} #1\par}
\vspace{0.18in}
{\Large\color{Teal} #2\par}
\vspace{0.5in}
\begin{tcolorbox}[colback=LightGray,colframe=Gray,boxrule=0.5pt,arc=2mm]
\textbf{Purpose.} #3
\end{tcolorbox}
\vfill
{\large Prepared from the revised edition of \textit{The Intelligent Investor}, with Jason Zweig's commentary.\par}
\vspace{0.12in}
{\small #4\par}
\vspace{0.25in}
{\small\color{Gray}\parbox{0.90\textwidth}{\centering Educational material only. Historical thresholds and examples are preserved for study; they are not individualized investment, tax, or legal advice.}}\par
\vspace{0.35in}
{\color{Navy}\rule{\linewidth}{2.2pt}}
\end{titlepage}
}

\newcommand{\DocShortTitle}{Technical Guide}
\begin{document}
\makedoccover{The Intelligent Investor - Technical Guide}{A source-grounded value-investing operating model}{Organizes the book into an analytical and behavioral framework suitable for repeated use.}{Standalone LaTeX study document}
\tableofcontents\newpage

\begin{cautionbox}[Source and modernization boundary]
The underlying source is the fourth revised edition, with Benjamin Graham's text updated in 1971--1972 and later commentary by Jason Zweig. Historical yields, tax rules, company examples, market levels, and dollar thresholds are therefore period-specific. This document preserves the source's framework and labels later implementation ideas separately rather than silently treating historical figures as current facts.
\end{cautionbox}

\section{Executive Framework}
\begin{principlebox}[The operating model]
The source treats successful investing less as a forecasting contest than as a disciplined operating system. Its recurring sequence is: distinguish investment from speculation; define a portfolio policy suited to temperament and effort; insist on price-conscious security selection; diversify; interpret market volatility as information rather than instruction; and preserve a margin of safety against analytical error and adverse outcomes.
\end{principlebox}

The book's central distinction is behavioral as much as analytical. An investment is evaluated through analysis, protection of principal, and an adequate return; activity that does not satisfy those conditions belongs in the speculative category. The practical implication is not that speculation can never occur, but that it should be recognized, sized, and segregated rather than disguised as investment.

\subsection{The six-part decision architecture}
\begin{enumerate}
\item \textbf{Classification:} Is the operation investment or speculation?
\item \textbf{Investor mode:} Is the owner defensive or enterprising based on willingness and ability to do sustained analytical work?
\item \textbf{Allocation:} What stock/bond mix can be maintained through adverse markets without forced behavior?
\item \textbf{Selection:} Does the security meet quality, financial-strength, earnings, and price tests?
\item \textbf{Behavior:} Can market quotations be treated as offers rather than commands?
\item \textbf{Margin of safety:} Is there enough valuation or coverage cushion to absorb error?
\end{enumerate}

\section{Investment Versus Speculation}
\subsection{Operational test}
The source's definition can be translated into an investment gate:
\begin{center}
\begin{tabularx}{\linewidth}{@{}p{0.24\linewidth}Y Y@{}}
\toprule
\textbf{Dimension} & \textbf{Investment standard} & \textbf{Speculative warning}\tabularnewline
\midrule
Analysis & Facts are studied against standards of safety and value. & Decision depends mainly on price action, narrative, or someone else paying more.\\
Principal & Serious permanent loss is deliberately guarded against. & Downside is undefined, ignored, or financed with leverage.\\
Return & Target is satisfactory and sustainable. & Outcome requires exceptional growth, perfect timing, or multiple expansion.\\
Price & Price is judged against business value. & Value is inferred from the quoted price itself.\\
Process & Rules remain stable across market moods. & Rules change after rallies, selloffs, headlines, or social proof.\\
\bottomrule
\end{tabularx}
\end{center}
\sourceRef{Introduction; Chapter 1 and commentary.}

\section{Defensive and Enterprising Investors}
The categories are not primarily labels for risk tolerance. They describe how much intelligent effort the investor can consistently devote to the task.

\begin{center}
\begin{tabularx}{\linewidth}{@{}p{0.22\linewidth}Y Y@{}}
\toprule
& \textbf{Defensive} & \textbf{Enterprising}\\
\midrule
Primary objective & Avoid serious mistakes and reduce decision burden. & Seek superior opportunities through disciplined additional work.\\
Portfolio style & Broad diversification, high-grade bonds, leading common stocks, funds/indexing. & Same sound base plus carefully justified departures.\\
Research burden & Low to moderate. & High and continuous.\\
Typical failure & Chasing performance or overpaying for glamour. & Confusing activity with skill; accepting weak quality for yield or excitement.\\
Best control & Simplicity, automation, low turnover. & Narrow mandate, explicit criteria, written thesis, post-analysis.\\
\bottomrule
\end{tabularx}
\end{center}

\section{Portfolio Policy}
\subsection{The Graham allocation band}
For the defensive investor, the text repeatedly uses a broad 25\%--75\% band for common stocks, with the inverse held in high-grade bonds and a 50/50 division as the standard reference point. The purpose is not tactical precision; it is to prevent extreme allocation decisions and create a rebalancing discipline.

\begin{equation}
0.25 \leq w_{equity} \leq 0.75, \qquad w_{bond}=1-w_{equity}
\end{equation}

A simple historical rule described in Chapter 1 is to restore the 50/50 balance when market movement produces a drift of roughly five percentage points. The enduring lesson is systematic rebalancing, not the historical number itself.

\subsection{Dollar-cost averaging}
Regular equal-dollar purchases are presented as a way to make the investor buy more shares at lower prices and fewer at higher prices without forecasting the market.
\begin{equation}
q_t=\frac{C}{P_t}, \qquad \bar{P}=\frac{\sum_t C_t}{\sum_t q_t}
\end{equation}
where $C_t$ is the contribution, $P_t$ the purchase price, and $q_t$ the shares acquired.

\section{Defensive Stock Selection}
Chapter 14 formalizes seven quality-and-price tests. Historical dollar cutoffs for company size are not portable, but the structure is.

\begin{center}
\small
\begin{tabularx}{\linewidth}{@{}p{0.06\linewidth}p{0.27\linewidth}Y@{}}
\toprule
\textbf{No.} & \textbf{Criterion} & \textbf{Source-era test / enduring intent}\\
\midrule
1 & Adequate size & Avoid fragile enterprises whose small scale adds unnecessary business risk.\\
2 & Strong financial condition & For industrials, current assets at least twice current liabilities; long-term debt not above net current assets.\\
3 & Earnings stability & Positive earnings in each of the past ten years.\\
4 & Dividend record & Long uninterrupted dividend history.\\
5 & Earnings growth & At least one-third increase in per-share earnings over ten years, using three-year averages at the endpoints.\\
6 & Moderate P/E & Price no more than 15 times average earnings of the past three years.\\
7 & Moderate price/book & Price no more than 1.5 times book value, with $P/E \times P/B \leq 22.5$ as a combined rule.\\
\bottomrule
\end{tabularx}
\end{center}

\begin{modernbox}[Derived shorthand]
The familiar ``Graham Number'' is a later algebraic shorthand derived from the combined Chapter 14 limits, not a separately named formula in the text:
\[
P_{max}\approx \sqrt{22.5\,EPS\,BVPS}.
\]
Use it only as a screening aid and only where book value is economically meaningful.
\end{modernbox}

\section{Enterprising Stock Selection}
The enterprising investor begins from the same sound base and earns the right to depart only through analysis. The text's positive approaches include unpopular large companies, bargain issues, special situations, and net-current-asset opportunities. Its negative approach rejects reaching for yield in weak bonds, chasing new issues, and buying complexity merely because it promises excitement.

\subsection{Net-current-asset approach}
The stripped-down asset value used in the classic bargain method is:
\begin{equation}
NCAV = Current\ Assets - Total\ Liabilities - Prior\ Claims.
\end{equation}
Graham describes historical purchases made typically at two-thirds or less of this value and emphasizes wide diversification. This is a specialized enterprising method, not a default recommendation for a defensive investor.

\subsection{Enterprising quantitative screen}
Chapter 15 illustrates a looser screen than the defensive standard: current assets at least 1.5 times current liabilities; debt not more than 110\% of net current assets for industrials; no deficit over five years; a current dividend; earnings improvement; and price below 120\% of net tangible assets. These are historical screening rules, not immutable contemporary thresholds.

\section{Security Analysis}
\subsection{Bond and preferred-stock safety}
Chapter 11 emphasizes coverage of fixed charges across a period of years, including a weak year, rather than a single favorable period. The idea is structural: fixed claims should be supported by earnings with room to deteriorate before principal or interest is threatened.

\subsection{Common-stock valuation}
The source warns that valuation becomes least reliable when analysts project high growth far into the future. One simplified historical formula is:
\begin{equation}
V = E(8.5 + 2g),
\end{equation}
where $E$ is current normal earnings and $g$ is the expected annual growth rate over roughly seven to ten years. The endnotes explicitly caution that this does not produce a ``true value''; it approximates more elaborate growth-stock calculations used at the time.

\subsection{Earnings quality}
Chapter 12 directs the reader away from overreacting to a single year's EPS and toward the accounting details beneath per-share figures: special charges, dilution, pro-forma adjustments, and other presentation choices. The durable procedure is to normalize earnings, read footnotes, compare accounting policies with peers, and treat unusually favorable adjustments skeptically.

\section{Market Fluctuations and Investor Psychology}
Chapter 8 separates \textbf{timing} from \textbf{pricing}. Timing attempts to predict the market's next direction. Pricing asks whether the security is cheap or dear relative to value. The book strongly favors the second discipline.

\begin{principlebox}[Mr. Market logic]
A quotation is an offer from a volatile counterparty. It may be useful when it offers an attractive transaction, but it is not an instruction about the underlying business. The investor's job is to exploit price/value gaps without allowing quotation changes to redefine value automatically.
\end{principlebox}

This is also the behavioral core of the book: the investor's own temperament can be a larger source of permanent loss than normal market volatility.

\section{Funds and Advisers}
Chapters 9 and 10 ask two different questions: whether an investor can identify superior funds in advance, and how to use professional advice without surrendering judgment. The commentary's practical defensive solution is low-cost diversified indexing. For advisers, the source emphasizes modest promises, conservative implementation, reasonable fees, and the investor's duty to verify trustworthiness.

\section{Corporate Management and Capital Allocation}
Chapters 18 and 19 shift attention from security selection to ownership. The shareholder should judge whether management uses retained earnings productively, pays sensible dividends, repurchases shares at rational prices, controls acquisition risk, and communicates as an owner-oriented steward. Capital allocation is treated as an investment problem inside the business.

\section{Case-History Failure Patterns}
Chapter 17's four cases create a reusable red-flag catalog:
\begin{itemize}
\item weak fixed-charge coverage and deteriorating operating economics;
\item debt-fueled empire building and acquisition cascades;
\item accounting complexity used to make a risky structure look stronger;
\item speculative financing of a small enterprise on the strength of a fashionable concept rather than demonstrated earning power.
\end{itemize}
The lesson is procedural: elementary balance-sheet and coverage tests should be completed before narratives, forecasts, or market stories are allowed to dominate the decision.

\section{Margin of Safety}
Chapter 20 identifies margin of safety as the organizing concept that ties the prior chapters together. For bonds, it appears as excess earnings over fixed charges. For common stocks, one expression is the excess of earnings yield over a high-grade bond yield, combined with diversification and a reasonable purchase price.

\begin{align}
Earnings\ Yield &= \frac{EPS}{Price}=\frac{1}{P/E},\\
MOS_{spread} &= Earnings\ Yield - Bond\ Yield.
\end{align}

\begin{cautionbox}[Do not reduce margin of safety to one formula]
The book uses margin of safety as a broad protection against error. It includes conservative assumptions, balance-sheet strength, price discipline, diversification, and the avoidance of low-quality securities sold under favorable business conditions. A single percentage cannot substitute for those controls.
\end{cautionbox}

\section{A Repeatable Investment Workflow}
\begin{enumerate}
\item Write the mandate: defensive or enterprising; time horizon; liquidity needs; maximum speculative allocation.
\item Set a stock/bond policy and a rebalancing rule before observing today's market mood.
\item Separate investing capital from speculation capital operationally and mentally.
\item For direct stocks, screen for balance-sheet strength, earnings durability, and valuation before narrative research.
\item Normalize EPS and review footnotes; reject accounting that cannot be understood.
\item Estimate value with simple, conservative methods; require a margin of safety.
\item Diversify so a single analytical error is survivable.
\item Record the thesis, disconfirming evidence, and exit conditions.
\item Review businesses on a schedule; do not monitor quotations as if they were business results.
\item Rebalance and continue contributions mechanically unless the underlying policy itself has changed.
\end{enumerate}

\section*{Primary Source}
Benjamin Graham, \textit{The Intelligent Investor: A Book of Practical Counsel}, revised edition, updated with commentary by Jason Zweig, preface and appendix material by Warren E. Buffett. Chapter references in this document point to the structure of that edition.

\end{document}
